\section*{5. 与线束相伴的态射 $\rho:\widetilde X\longrightarrow\mathbb P^1$ 的上同调研究}

\subsection*{5.1. 限制同态
$H^*(\widetilde X)\longrightarrow H^*(Y)$}

以 $Y$ 表示 $\rho:\widetilde X\to\mathbb P^1$ 的一个光滑纤维, 设为
\[
Y=\rho^{-1}(t_0)=X\mathbin{\cdot}H_{t_0}.
\]
有交换图
\begin{center}
\resizebox{\textwidth}{!}{%
\begin{tikzcd}[ampersand replacement=\&,column sep=large,row sep=huge]
\mathbb P^1
  \& {} \& \widetilde X\arrow[ll,"\rho"']\arrow[d,"f"]
  \& {} \& \widetilde\Delta
       \arrow[ll,hook',"i"]\arrow[d,"g"]
       \arrow[r,phantom,"\simeq"]
  \& \Delta\times\mathbb P^1 \\
\mathrm{Sp}(k)=\{t_0\}\arrow[u,hook]
  \& {} \& X
  \& {} \& \Delta
       \arrow[ll,hook',"j"]
       \arrow[ur,hook,"{\ell=\mathrm{id}\times\{t_0\}}"']
       \arrow[dlll,hook,"h"]
  \& {} \\
{} \& Y\arrow[ul]\arrow[uur,"k"]\arrow[ur,"m"']
  \& {} \& {} \& {} \& {}
\end{tikzcd}}
\end{center}

\subsection*{命题 5.1.1.}
\underline{在 (3.1) 的假设下, 以 $Y$ 表示一个光滑纤维:}

\smallskip
\noindent
\textup{(5.1.2)} \underline{限制同态}
\[
k^*:H^\bullet(\widetilde X,\mathbb Z_\ell)
   \longrightarrow H^\bullet(Y,\mathbb Z_\ell)
\]
\underline{经由 (4.2.2) 表示为}
\[
H^\bullet(X,\mathbb Z_\ell)\oplus
H^{\bullet-2}(\Delta,\mathbb Z_\ell(-1))
\xrightarrow{\ m^*+h_*\ }
H^\bullet(Y,\mathbb Z_\ell).
\]

\noindent
\textup{(5.1.3)} \underline{Gysin 映射}
\[
k_*:H^{\bullet-2}(Y,\mathbb Z_\ell(-1))
   \longrightarrow H^\bullet(\widetilde X,\mathbb Z_\ell)
\]
\underline{表示为}
\[
H^{\bullet-2}(Y,\mathbb Z_\ell(-1))
\xrightarrow{\ m_*\oplus(-h^*)\ }
H^\bullet(X,\mathbb Z_\ell)\oplus
H^{\bullet-2}(\Delta,\mathbb Z_\ell(-1)).
\]

\underline{证明.}
$\widetilde X$ 的光滑子概形 $\widetilde\Delta$ 和 $Y$ 横截相交于 $\Delta$; 换言之, 光滑概形的方块
\begin{equation}\tag{5.1.4}
\begin{tikzcd}[column sep=large,row sep=large]
\Delta\arrow[r,"\ell"]\arrow[d,hook,"h"]
  &\widetilde\Delta\arrow[d,hook,"i"]\\
Y\arrow[r,"k"']&\widetilde X
\end{tikzcd}
\end{equation}
为 笛卡儿方块. 因此由它可得公式 [SGA 5, IV]
\begin{equation}\tag{5.1.5}
h_*\ell^*=k^*i_* ,\qquad \ell_*h^*=i^*k_*.
\end{equation}
复合
\[
H^\bullet(X,\mathbb Z_\ell)\oplus
H^{\bullet-2}(\Delta,\mathbb Z_\ell(-1))
\xrightarrow{\ f^*+i_*g^*\ }
H^\bullet(\widetilde X,\mathbb Z_\ell)
\xrightarrow{\ k^*\ }
H^\bullet(Y,\mathbb Z_\ell)
\]
为 $k^*f^*=m^*$ 与 $k^*i_*g^*=h_*\ell^*g^*$ 之和, 后一个等式由 (5.1.5) 得到. 由于 $g\ell=\mathrm{id}_\Delta$, 这就证明了 (5.1.2).

复合
\[
H^{\bullet-2}(Y,\mathbb Z_\ell(-1))
\xrightarrow{\ k_*\ }
H^\bullet(\widetilde X,\mathbb Z_\ell)
\xrightarrow{\ f_*\oplus(-g_*i^*)\ }
H^\bullet(X,\mathbb Z_\ell)\oplus
H^{\bullet-2}(\Delta,\mathbb Z_\ell(-1))
\]
为下列两式之和
\[
f_*k_*=m_*
\qquad\text{及}\qquad
-g_*i^*k_*=-g_*\ell_*h^*=-h^*,
\]
这就证明了 (5.1.3).

\subsection*{推论 5.1.6.}
\underline{在 5.1.1 的假设下, 当 $i\leq n-1$ 时, 映射}
\[
m^*:H^i(X,\mathbb Z_\ell)\longrightarrow H^i(Y,\mathbb Z_\ell)
\]
\underline{与映射}
\[
k^*:H^i(\widetilde X,\mathbb Z_\ell)
   \longrightarrow H^i(Y,\mathbb Z_\ell)
\]
\underline{的像相等.}

\underline{证明.}
由超平面截面的弱 Lefschetz 定理 (SGA 5, VII 7.1), 当 $i\leq n-1$ 时, 映射
\begin{equation}\tag{5.1.7}
H^{i-2}(X,\mathbb Z_\ell(-1))
\xrightarrow{\ j^*\ }
H^{i-2}(\Delta,\mathbb Z_\ell(-1))
\end{equation}
为同构. 因而由 (5.1.2), $k^*$ 的像为 $m^*$ 的像与下列复合的像之和:
\[
H^{i-2}(X,\mathbb Z_\ell(-1))
\xrightarrow{\ j^*\ }
H^{i-2}(\Delta,\mathbb Z_\ell(-1))
\xrightarrow{\ h_*\ }
H^i(Y,\mathbb Z_\ell).
\]
但投影公式给出
\[
\begin{aligned}
h_*j^*(x)
 &=h_*\bigl(h^*m^*(x)\bigr)
  =m^*(x)h_*(1)
  =\xi_Y\,m^*(x)\\
 &=m^*(x\xi_X),
 \qquad\text{因为 }m^*(\xi_X)=\xi_Y.
\end{aligned}
\]
因而 $\operatorname{Im}k^*\subset\operatorname{Im}m^*$, 而反向包含直接由 (5.1.2) 得出. 推论得证.
